\textbf{Phần 1. Trắc nghiệm lựa chọn}
\setcounterpageref{ex}{0}

\begin{ex} %[2H3V16]
Trong can có $15\text{ L}$ dầu. Người ta rót vào $3$ chai, mỗi chai $2\text{ L}$. Hỏi trong can còn lại bao nhiêu lít dầu?
\choice
{$10\text{ L}$}
{$8\text{ L}$}
{$11\text{ L}$}
{\True $9\text{ L}$}
\loigiai{
Số dầu đã rót: $2 \times 3 = 6\text{ (L)}$. Trong can còn lại: $15 - 6 = 9\text{ (L)}$.
\textbf{Đáp án của câu hỏi là D}
}
\end{ex}

\begin{ex} %[2H3V16]
Thùng thứ nhất chứa $38\text{ L}$ nước, thùng thứ hai chứa ít hơn thùng thứ nhất $9\text{ L}$ nước. Cả hai thùng chứa bao nhiêu lít nước?
\choice
{$65\text{ L}$}
{\True $67\text{ L}$}
{$69\text{ L}$}
{$63\text{ L}$}
\loigiai{
Thùng thứ hai: $38 - 9 = 29\text{ (L)}$.\nCả hai thùng: $38 + 29 = 67\text{ (L)}$.
\textbf{Đáp án của câu hỏi là B}
}
\end{ex}

\begin{ex} %[2H3V16]
Tính giá trị: $45\text{ L} + 27\text{ L} - 18\text{ L}$.
\choice
{$52\text{ L}$}
{$56\text{ L}$}
{\True $54\text{ L}$}
{$50\text{ L}$}
\loigiai{
$45 + 27 = 72\text{ L}$; $72 - 18 = 54\text{ L}$.
\textbf{Đáp án của câu hỏi là C}
}
\end{ex}

\begin{ex} %[2H3V16]
Một bể chứa $80\text{ L}$ nước. Buổi sáng người ta dùng $25\text{ L}$, buổi chiều dùng $36\text{ L}$. Hỏi trong bể còn lại bao nhiêu lít nước?
\choice
{\True $19\text{ L}$}
{$21\text{ L}$}
{$18\text{ L}$}
{$23\text{ L}$}
\loigiai{
Tổng số lít nước đã dùng: $25 + 36 = 61\text{ (L)}$.\nTrong bể còn lại: $80 - 61 = 19\text{ (L)}$.
\textbf{Đáp án của câu hỏi là A}
}
\end{ex}

\begin{ex} %[2H3V16]
Có hai can nước mắm: can to chứa $20\text{ L}$, can nhỏ chứa ít hơn can to $8\text{ L}$. Cả hai can chứa bao nhiêu lít nước mắm?
\choice
{$30\text{ L}$}
{$28\text{ L}$}
{$34\text{ L}$}
{\True $32\text{ L}$}
\loigiai{
Can nhỏ chứa: $20 - 8 = 12\text{ (L)}$.\nCả hai can chứa: $20 + 12 = 32\text{ (L)}$.
\textbf{Đáp án của câu hỏi là D}
}
\end{ex}

\begin{ex} %[2H3V16]
Rót hết nước từ một bình $10\text{ L}$ vào các ca loại $2\text{ L}$. Hỏi rót được bao nhiêu ca nước?
\choice
{$4$ ca}
{\True $5$ ca}
{$6$ ca}
{$8$ ca}
\loigiai{
$10 : 2 = 5$ (ca).
\textbf{Đáp án của câu hỏi là B}
}
\end{ex}

\begin{ex} %[2H3V16]
Điền số thích hợp: $18\text{ L} + \dots\text{ L} = 35\text{ L}$.
\choice
{$16$}
{$18$}
{\True $17$}
{$15$}
\loigiai{
$35 - 18 = 17$.
\textbf{Đáp án của câu hỏi là C}
}
\end{ex}

\begin{ex} %[2H3V16]
Mẹ mua về $12\text{ L}$ sữa tươi. Mẹ chia đều cho $4$ chị em uống trong một tuần. Hỏi mỗi người được uống bao nhiêu lít sữa?
\choice
{\True $3\text{ L}$}
{$2\text{ L}$}
{$4\text{ L}$}
{$5\text{ L}$}
\loigiai{
$12 : 4 = 3\text{ (L)}$.
\textbf{Đáp án của câu hỏi là A}
}
\end{ex}

\begin{ex} %[2H3V16]
Một cửa hàng xăng dầu có $95\text{ L}$ dầu. Ngày thứ nhất bán được $37\text{ L}$, ngày thứ hai bán được $29\text{ L}$. Sau hai ngày cửa hàng còn lại bao nhiêu lít dầu?
\choice
{$31\text{ L}$}
{$27\text{ L}$}
{$33\text{ L}$}
{\True $29\text{ L}$}
\loigiai{
Số dầu đã bán: $37 + 29 = 66\text{ (L)}$.\nCòn lại: $95 - 66 = 29\text{ (L)}$.
\textbf{Đáp án của câu hỏi là D}
}
\end{ex}

\begin{ex} %[2H3V16]
Can thứ nhất có $14\text{ L}$ mật ong. Nếu chuyển $3\text{ L}$ từ can thứ nhất sang can thứ hai thì hai can có số lít mật ong bằng nhau. Hỏi can thứ hai lúc đầu có bao nhiêu lít mật ong?
\choice
{$9\text{ L}$}
{\True $8\text{ L}$}
{$11\text{ L}$}
{$7\text{ L}$}
\loigiai{
Sau khi chuyển, mỗi can có: $14 - 3 = 11\text{ (L)}$.\nLúc đầu can thứ hai có: $11 - 3 = 8\text{ (L)}$.
\textbf{Đáp án của câu hỏi là B}
}
\end{ex}

\begin{ex} %[2H3V16]
Người ta đong được $3$ can $5\text{ L}$ và $1$ can $2\text{ L}$. Tổng số lít nước đong được là
\choice
{$15\text{ L}$}
{$18\text{ L}$}
{\True $17\text{ L}$}
{$16\text{ L}$}
\loigiai{
$5 \times 3 + 2 = 15 + 2 = 17\text{ (L)}$.
\textbf{Đáp án của câu hỏi là C}
}
\end{ex}

\begin{ex} %[2H3V16]
Hiệu giữa số lít nước của can $25\text{ L}$ và can $18\text{ L}$ là
\choice
{\True $7\text{ L}$}
{$6\text{ L}$}
{$8\text{ L}$}
{$9\text{ L}$}
\loigiai{
$25 - 18 = 7\text{ (L)}$.
\textbf{Đáp án của câu hỏi là A}
}
\end{ex}

\textbf{Phần 2. Trắc nghiệm đúng sai}
\setcounterpageref{ex}{0}

\begin{ex} %[2H3V16]
Xét tính đúng sai của các nhận định về đơn vị đo dung tích (Lít):
\choiceTF
{\True Lít là đơn vị dùng để đo sức chứa của một vật (dung tích), viết tắt là L hoặc l.}
{Can 10 lít chứa được lượng nước ít hơn 2 can 5 lít cộng lại.}
{\True Muốn đong được đúng 3 lít nước bằng can 5 lít và can 2 lít, ta múc đầy can 5 lít rồi rót đầy vào can 2 lít, số nước còn lại trong can 5 lít đúng bằng 3 lít.}
{\True Một bể nước có thể tích 50 L, đã có 35 L nước, cần đổ thêm 15 L để đầy bể.}
\loigiai{
\begin{itemize}
    \item Ý a) Đúng: Chuẩn kiến thức sách giáo khoa.
    \item Ý b) Sai: $5 + 5 = 10\text{ L}$, hai bên bằng nhau.
    \item Ý c) Đúng: $5 - 2 = 3\text{ L}$. Đúng.
    \item Ý d) Đúng: $50 - 35 = 15\text{ L}$. Đúng.
\end{itemize}
\begin{itemize}
    \item a) \textbf{Đúng}.
    \item b) \textbf{Sai}.
    \item c) \textbf{Đúng}.
    \item d) \textbf{Đúng}.
\end{itemize}
}
\end{ex}

\begin{ex} %[2H3V16]
Có ba thùng đựng dầu: Thùng A có 25 L, Thùng B có 35 L, Thùng C có 40 L:
\choiceTF
{\True Cả ba thùng có tất cả 100 L dầu.}
{\True Thùng C nhiều hơn thùng A là 15 L dầu.}
{Thùng B nhiều hơn thùng C là 5 L dầu.}
{Chuyển 5 L từ thùng C sang thùng A thì ba thùng có số dầu bằng nhau.}
\loigiai{
\begin{itemize}
    \item Ý a) Đúng: $25 + 35 + 40 = 100\text{ L}$.
    \item Ý b) Đúng: $40 - 25 = 15\text{ L}$.
    \item Ý c) Sai: Thùng B (35 L) ít hơn Thùng C (40 L) là 5 L.
    \item Ý d) Sai: Khi đó A có 30 L, B có 35 L, C có 35 L, chưa bằng nhau.
\end{itemize}
\begin{itemize}
    \item a) \textbf{Đúng}.
    \item b) \textbf{Đúng}.
    \item c) \textbf{Sai}.
    \item d) \textbf{Sai}.
\end{itemize}
}
\end{ex}

\textbf{Phần 3. Trả lời ngắn}
\setcounterpageref{ex}{0}

\begin{ex} %[2H3V16]
Một can dầu chứa 18 lít dầu. Bác Nam rót ra 2 can, mỗi can 5 lít. Hỏi trong can ban đầu còn lại bao nhiêu lít dầu?
\shortans{8}
\loigiai{
\textbf{Đáp số:} $8$\\
Số dầu rót ra: $5 + 5 = 10\text{ L}$. Còn lại: $18 - 10 = 8\text{ L}$.
}
\end{ex}

\begin{ex} %[2H3V16]
Thùng thứ nhất có 36 lít nước, thùng thứ hai có ít hơn thùng thứ nhất 8 lít nước. Hỏi thùng thứ hai có bao nhiêu lít nước?
\shortans{28}
\loigiai{
\textbf{Đáp số:} $28$\\
$36 - 8 = 28\text{ L}$.
}
\end{ex}

\begin{ex} %[2H3V16]
Tính giá trị biểu thức: $42\text{ L} - 15\text{ L} + 18\text{ L}$.
\shortans{45}
\loigiai{
\textbf{Đáp số:} $45$\\
$42 - 15 = 27$; $27 + 18 = 45$.
}
\end{ex}

\begin{ex} %[2H3V16]
Người ta dùng các can loại 5 lít để đong hết 20 lít dầu. Hỏi cần bao nhiêu chiếc can loại đó?
\shortans{4}
\loigiai{
\textbf{Đáp số:} $4$\\
$20 : 5 = 4$ chiếc can.
}
\end{ex}

\begin{ex} %[2H3V16]
Một bể nước chứa 65 lít nước, người ta xả bớt đi 28 lít nước. Hỏi trong bể còn lại bao nhiêu lít nước?
\shortans{37}
\loigiai{
\textbf{Đáp số:} $37$\\
$65 - 28 = 37\text{ L}$.
}
\end{ex}

\begin{ex} %[2H3V16]
Can A có 15 lít nước mắm, Can B có nhiều hơn Can A 7 lít nước mắm. Hỏi cả hai can có bao nhiêu lít nước mắm?
\shortans{37}
\loigiai{
\textbf{Đáp số:} $37$\\
Can B: $15 + 7 = 22\text{ L}$. Cả hai can: $15 + 22 = 37\text{ L}$.
}
\end{ex}
